\documentclass[10pt,a4paper]{amsart}
\usepackage{amsmath,amssymb,amsthm,graphicx}
% Retained minimal preamble for the Stacks Project chapter flat.tex (More on Flatness).
% Verbatim excerpt of the chapter preamble: Xy-pic setup (native preamble lines 14-18) and the native
% theorem declarations with the mathematical macros (native preamble lines 154-217). Copied mechanically;
% nothing added, removed or redefined. The wrapper class and the amsmath/amssymb/amsthm/graphicx base
% packages are supplied by the assembler, so the class/externaldocument part of the native preamble is omitted.
% For commutative diagrams we use Xy-pic
\usepackage[all]{xy}

% We use 2cell for 2-commutative diagrams.
\xyoption{2cell}
% Theorem environments.
%
\theoremstyle{plain}
\newtheorem{theorem}[subsection]{Theorem}
\newtheorem{proposition}[subsection]{Proposition}
\newtheorem{lemma}[subsection]{Lemma}

\theoremstyle{definition}
\newtheorem{definition}[subsection]{Definition}
\newtheorem{example}[subsection]{Example}
\newtheorem{exercise}[subsection]{Exercise}
\newtheorem{situation}[subsection]{Situation}

\theoremstyle{remark}
\newtheorem{remark}[subsection]{Remark}
\newtheorem{remarks}[subsection]{Remarks}

\numberwithin{equation}{subsection}

% Macros
%
\def\lim{\mathop{\mathrm{lim}}\nolimits}
\def\colim{\mathop{\mathrm{colim}}\nolimits}
\def\Spec{\mathop{\mathrm{Spec}}}
\def\Hom{\mathop{\mathrm{Hom}}\nolimits}
\def\Ext{\mathop{\mathrm{Ext}}\nolimits}
\def\SheafHom{\mathop{\mathcal{H}\!\mathit{om}}\nolimits}
\def\SheafExt{\mathop{\mathcal{E}\!\mathit{xt}}\nolimits}
\def\Sch{\mathit{Sch}}
\def\Mor{\mathop{\mathrm{Mor}}\nolimits}
\def\Ob{\mathop{\mathrm{Ob}}\nolimits}
\def\Sh{\mathop{\mathit{Sh}}\nolimits}
\def\NL{\mathop{N\!L}\nolimits}
\def\CH{\mathop{\mathrm{CH}}\nolimits}
\def\proetale{{pro\text{-}\acute{e}tale}}
\def\etale{{\acute{e}tale}}
\def\QCoh{\mathit{QCoh}}
\def\Ker{\mathop{\mathrm{Ker}}}
\def\Im{\mathop{\mathrm{Im}}}
\def\Coker{\mathop{\mathrm{Coker}}}
\def\Coim{\mathop{\mathrm{Coim}}}

% Boxtimes
%
\DeclareMathSymbol{\boxtimes}{\mathbin}{AMSa}{"02}

%
% Macros for moduli stacks/spaces
%
\def\QCohstack{\mathcal{QC}\!\mathit{oh}}
\def\Cohstack{\mathcal{C}\!\mathit{oh}}
\def\Spacesstack{\mathcal{S}\!\mathit{paces}}
\def\Quotfunctor{\mathrm{Quot}}
\def\Hilbfunctor{\mathrm{Hilb}}
\def\Curvesstack{\mathcal{C}\!\mathit{urves}}
\def\Polarizedstack{\mathcal{P}\!\mathit{olarized}}
\def\Complexesstack{\mathcal{C}\!\mathit{omplexes}}
% \Pic is the operator that assigns to X its picard group, usage \Pic(X)
% \Picardstack_{X/B} denotes the Picard stack of X over B
% \Picardfunctor_{X/B} denotes the Picard functor of X over B
\def\Pic{\mathop{\mathrm{Pic}}\nolimits}
\def\Picardstack{\mathcal{P}\!\mathit{ic}}
\def\Picardfunctor{\mathrm{Pic}}
\def\Deformationcategory{\mathcal{D}\!\mathit{ef}}

% --- appended 2026-10-12 (Chief): the project own definitions of the reference/slogan/history
% display environments, copied mechanically from the shared Stacks preamble lines 8-12
% (\usepackage{verbatim} plus the three \newenvironment lines, which the chapter preamble keeps
% verbatim from the project preamble). Nothing added, removed or redefined: these environments
% carry no mathematical content and the project itself typesets them as comments. Without these
% three lines a counted statement whose own text uses them fails with Environment slogan undefined.
\usepackage{verbatim}
\newenvironment{reference}{\comment}{\endcomment}
\newenvironment{slogan}{\comment}{\endcomment}
\newenvironment{history}{\comment}{\endcomment}
\usepackage{cleveref}
\crefname{originalsection}{original section}{original sections}
\Crefname{originalsection}{Original section}{Original sections}
\makeatletter
\newlabel{schemes-section-quasi-compact}{{01K2}{source}{}{sourcepointer-1}{}}
\newlabel{schemes-section-quasi-compact@cref}{{[originalsection][0][]01K2}{[1][0][]source}{}{}{}}
\makeatother
\crefname{originalsection}{original section}{original sections}
\Crefname{originalsection}{Original section}{Original sections}
\makeatletter
\newlabel{schemes-section-separation-axioms}{{01KH}{source}{}{sourcepointer-2}{}}
\newlabel{schemes-section-separation-axioms@cref}{{[originalsection][0][]01KH}{[1][0][]source}{}{}{}}
\makeatother
\crefname{originallemma}{original lemma}{original lemmas}
\Crefname{originallemma}{Original lemma}{Original lemmas}
\makeatletter
\newlabel{morphisms-lemma-quasi-compact-immersion}{{01RG}{source}{}{sourcepointer-3}{}}
\newlabel{morphisms-lemma-quasi-compact-immersion@cref}{{[originallemma][0][]01RG}{[1][0][]source}{}{}{}}
\makeatother
\crefname{originallemma}{original lemma}{original lemmas}
\Crefname{originallemma}{Original lemma}{Original lemmas}
\makeatletter
\newlabel{morphisms-lemma-scheme-theoretic-image-of-partial-section}{{0CNG}{source}{}{sourcepointer-4}{}}
\newlabel{morphisms-lemma-scheme-theoretic-image-of-partial-section@cref}{{[originallemma][0][]0CNG}{[1][0][]source}{}{}{}}
\makeatother
\crefname{originallemma}{original lemma}{original lemmas}
\Crefname{originallemma}{Original lemma}{Original lemmas}
\makeatletter
\newlabel{lemma-zariski-after-blowup}{{081S}{source}{}{sourcepointer-5}{}}
\newlabel{lemma-zariski-after-blowup@cref}{{[originallemma][0][]081S}{[1][0][]source}{}{}{}}
\makeatother
\crefname{originallemma}{original lemma}{original lemmas}
\Crefname{originallemma}{Original lemma}{Original lemmas}
\makeatletter
\newlabel{divisors-lemma-strict-transform}{{080E}{source}{}{sourcepointer-6}{}}
\newlabel{divisors-lemma-strict-transform@cref}{{[originallemma][0][]080E}{[1][0][]source}{}{}{}}
\makeatother
\crefname{originallemma}{original lemma}{original lemmas}
\Crefname{originallemma}{Original lemma}{Original lemmas}
\makeatletter
\newlabel{divisors-lemma-blowing-up-two-ideals}{{080A}{source}{}{sourcepointer-7}{}}
\newlabel{divisors-lemma-blowing-up-two-ideals@cref}{{[originallemma][0][]080A}{[1][0][]source}{}{}{}}
\makeatother
\crefname{originallemma}{original lemma}{original lemmas}
\Crefname{originallemma}{Original lemma}{Original lemmas}
\makeatletter
\newlabel{divisors-lemma-extend-admissible-blowups}{{080M}{source}{}{sourcepointer-8}{}}
\newlabel{divisors-lemma-extend-admissible-blowups@cref}{{[originallemma][0][]080M}{[1][0][]source}{}{}{}}
\makeatother
\crefname{originallemma}{original lemma}{original lemmas}
\Crefname{originallemma}{Original lemma}{Original lemmas}
\makeatletter
\newlabel{morphisms-lemma-image-proper-scheme-closed}{{01W6}{source}{}{sourcepointer-9}{}}
\newlabel{morphisms-lemma-image-proper-scheme-closed@cref}{{[originallemma][0][]01W6}{[1][0][]source}{}{}{}}
\makeatother
\crefname{originallemma}{original lemma}{original lemmas}
\Crefname{originallemma}{Original lemma}{Original lemmas}
\makeatletter
\newlabel{lemma-check-separated}{{0F3U}{source}{}{sourcepointer-10}{}}
\newlabel{lemma-check-separated@cref}{{[originallemma][0][]0F3U}{[1][0][]source}{}{}{}}
\makeatother
\crefname{originallemma}{original lemma}{original lemmas}
\Crefname{originallemma}{Original lemma}{Original lemmas}
\makeatletter
\newlabel{divisors-lemma-composition-admissible-blowups}{{080L}{source}{}{sourcepointer-11}{}}
\newlabel{divisors-lemma-composition-admissible-blowups@cref}{{[originallemma][0][]080L}{[1][0][]source}{}{}{}}
\makeatother
\title{Two open immersions of finite type separated schemes over a common open can be made part of one separated finite type scheme by admissible blowups}
\author{The Stacks Project Authors}
\date{stacks-project:flat.tex:local-snapshot-20261011; GNU FDL 1.2 or later}
\begin{document}
\maketitle
\noindent\textit{Editorial note.} All delivered mathematical text is the author's own, copied from the pinned Stacks Project chapter file flat.tex; no statement or proof step was rewritten, reordered or removed.\par\medskip
\noindent\textit{Editorial note.} The preamble is the chapter's own retained minimal preamble (sources/stacks-flat/retained-preamble.tex), a verbatim excerpt of the native Xy-pic setup and of the native theorem declarations and mathematical macros, followed by the project own definitions of the reference, slogan and history display environments (verbatim package plus three newenvironment lines) that the retained file carries at its end and that the project itself typesets as comments. No macro was invented or redefined, and the wrapper class and the amsmath/amssymb/amsthm/graphicx base packages are supplied by the assembler.\par\medskip
\noindent\textit{Editorial note.} The counted statement and proof contain no cross reference, no citation and no included file, so the delivered body is self contained.\par\medskip
\noindent\textit{Editorial note.} Canonical location of the counted proposition: Stacks tag 0F3X (sources/stacks-flat/pinned-original-source.tex line 9647).\par\medskip
\noindent\textit{Allowed standard prerequisites (editorial).}\begin{itemize}
\item Scheme theoretic images of immersions and partial sections. Morphisms, Lemmas labelled morphisms-lemma-quasi-compact-immersion and morphisms-lemma-scheme-theoretic-image-of-partial-section, cited at proof lines 9670-9679. Produces X\_12 as a closed subscheme of X\_1 x\_S X\_2 with the projections isomorphisms over U.
\item Making a scheme theoretic image an open subscheme by an admissible blowup. Same chapter, Lemma labelled lemma-zariski-after-blowup, cited at proof lines 9682-9686; Divisors, Lemmas labelled divisors-lemma-strict-transform, divisors-lemma-blowing-up-two-ideals and divisors-lemma-extend-admissible-blowups, cited at proof lines 9690-9706. Supplies the blowups making the strict transforms open and the blowup of X\_12 dominating both.
\item Properness of the immersion of the strict transform and closedness of proper images. Morphisms, Lemma labelled morphisms-lemma-image-proper-scheme-closed, cited at proof lines 9702-9706. Concludes that X'\_12 -> X\_1' x\_S X\_2' is a closed immersion, so the glueing is legitimate.
\item Separatedness of the glued scheme and composition of admissible blowups. Same chapter, Lemma labelled lemma-check-separated, cited at proof lines 9718-9720; Divisors, Lemma labelled divisors-lemma-composition-admissible-blowups, cited at proof lines 9722-9724. Gives separatedness of X over S and closes the proof.
\end{itemize}\medskip
\ifdefined\hypertarget\hypertarget{sourcepointer-1}{}\fi
\noindent\textit{Source reference.} Original section 01K2: sources/stacks-chow/tags.txt, source lines 2035--2035. Cross-chapter standard prerequisite located in the project tag list staged in the local snapshot; invoked at proof line 9667 of the counted proof. It is retained only as a navigation pointer; every load-bearing step is inside the delivered span.\par\medskip
\ifdefined\hypertarget\hypertarget{sourcepointer-2}{}\fi
\noindent\textit{Source reference.} Original section 01KH: sources/stacks-chow/tags.txt, source lines 2050--2050. Cross-chapter standard prerequisite located in the project tag list staged in the local snapshot; invoked at proof line 9668 of the counted proof. It is retained only as a navigation pointer; every load-bearing step is inside the delivered span.\par\medskip
\ifdefined\hypertarget\hypertarget{sourcepointer-3}{}\fi
\noindent\textit{Source reference.} Original lemma 01RG: sources/stacks-chow/tags.txt, source lines 2301--2301. Cross-chapter standard prerequisite located in the project tag list staged in the local snapshot; invoked at proof line 9674 of the counted proof. It is retained only as a navigation pointer; every load-bearing step is inside the delivered span.\par\medskip
\ifdefined\hypertarget\hypertarget{sourcepointer-4}{}\fi
\noindent\textit{Source reference.} Original lemma 0CNG: sources/stacks-chow/tags.txt, source lines 15891--15891. Cross-chapter standard prerequisite located in the project tag list staged in the local snapshot; invoked at proof line 9681 of the counted proof. It is retained only as a navigation pointer; every load-bearing step is inside the delivered span.\par\medskip
\ifdefined\hypertarget\hypertarget{sourcepointer-5}{}\fi
\noindent\textit{Source reference.} Original lemma 081S: sources/stacks-flat/pinned-original-source.tex, source lines 9197--9197. Same-chapter pointer into the author environment lemma at source line 9197; invoked at proof line 9685 of the counted proof. It is retained only as a navigation pointer; every load-bearing step is inside the delivered span.\par\medskip
\ifdefined\hypertarget\hypertarget{sourcepointer-6}{}\fi
\noindent\textit{Source reference.} Original lemma 080E: sources/stacks-chow/tags.txt, source lines 10179--10179. Cross-chapter standard prerequisite located in the project tag list staged in the local snapshot; invoked at proof line 9690 of the counted proof. It is retained only as a navigation pointer; every load-bearing step is inside the delivered span.\par\medskip
\ifdefined\hypertarget\hypertarget{sourcepointer-7}{}\fi
\noindent\textit{Source reference.} Original lemma 080A: sources/stacks-chow/tags.txt, source lines 10175--10175. Cross-chapter standard prerequisite located in the project tag list staged in the local snapshot; invoked at proof line 9693 of the counted proof. It is retained only as a navigation pointer; every load-bearing step is inside the delivered span.\par\medskip
\ifdefined\hypertarget\hypertarget{sourcepointer-8}{}\fi
\noindent\textit{Source reference.} Original lemma 080M: sources/stacks-chow/tags.txt, source lines 10187--10187. Cross-chapter standard prerequisite located in the project tag list staged in the local snapshot; invoked at proof line 9704 of the counted proof. It is retained only as a navigation pointer; every load-bearing step is inside the delivered span.\par\medskip
\ifdefined\hypertarget\hypertarget{sourcepointer-9}{}\fi
\noindent\textit{Source reference.} Original lemma 01W6: sources/stacks-chow/tags.txt, source lines 2471--2471. Cross-chapter standard prerequisite located in the project tag list staged in the local snapshot; invoked at proof line 9717 of the counted proof. It is retained only as a navigation pointer; every load-bearing step is inside the delivered span.\par\medskip
\ifdefined\hypertarget\hypertarget{sourcepointer-10}{}\fi
\noindent\textit{Source reference.} Original lemma 0F3U: sources/stacks-flat/pinned-original-source.tex, source lines 9540--9540. Same-chapter pointer into the author environment lemma at source line 9540; invoked at proof line 9721 of the counted proof. It is retained only as a navigation pointer; every load-bearing step is inside the delivered span.\par\medskip
\ifdefined\hypertarget\hypertarget{sourcepointer-11}{}\fi
\noindent\textit{Source reference.} Original lemma 080L: sources/stacks-chow/tags.txt, source lines 10186--10186. Cross-chapter standard prerequisite located in the project tag list staged in the local snapshot; invoked at proof line 9723 of the counted proof. It is retained only as a navigation pointer; every load-bearing step is inside the delivered span.\par\medskip
\begin{lemma}
\label{lemma-find-common-blowups}
Let $S$ be a quasi-compact and quasi-separated scheme.
Let $U \to X_1$ and $U \to X_2$ be open immersions
of schemes over $S$ and assume $U$, $X_1$, $X_2$ of finite
type and separated over $S$. Then there exists a commutative diagram
$$
\xymatrix{
X_1' \ar[d] \ar[r] & X & X_2' \ar[l] \ar[d] \\
X_1 & U \ar[l] \ar[lu] \ar[u] \ar[ru] \ar[r] & X_2
}
$$
of schemes over $S$ where $X_i' \to X_i$ is a $U$-admissible
blowup, $X_i' \to X$ is an open immersion, and $X$ is separated and finite
type over $S$.
\end{lemma}

\begin{proof}
Throughout the proof all schemes will be separated of finite type over $S$.
This in particular implies these schemes are quasi-compact and quasi-separated
and the morphisms between them are quasi-compact and separated.
See Schemes, Sections \ref{schemes-section-quasi-compact} and
\ref{schemes-section-separation-axioms}.
We will use that if $U \to W$ is an immersion of such schemes over $S$,
then the scheme theoretic image $Z$ of $U$ in $W$ is a closed subscheme
of $W$ and $U \to Z$ is an open immersion, $U \subset Z$ is scheme
theoretically dense, and $U \subset Z$ is dense topologically. See
Morphisms, Lemma
\ref{morphisms-lemma-quasi-compact-immersion}.

\medskip\noindent
Let $X_{12} \subset X_1 \times_S X_2$ be the scheme theoretic image
of $U \to X_1 \times_S X_2$. The projections $p_i : X_{12} \to X_i$
induce isomorphisms $p_i^{-1}(U) \to U$ by
Morphisms, Lemma
\ref{morphisms-lemma-scheme-theoretic-image-of-partial-section}.
Choose a $U$-admissible blowup $X_i^i \to X_i$ such that
the strict transform $X_{12}^i$ of $X_{12}$ is isomorphic to an
open subscheme of $X_i^i$, see
Lemma \ref{lemma-zariski-after-blowup}.
Let $\mathcal{I}_i \subset \mathcal{O}_{X_i}$ be the corresponding
finite type quasi-coherent sheaf of ideals.
Recall that $X_{12}^i \to X_{12}$ is the blowup in
$p_i^{-1}\mathcal{I}_i \mathcal{O}_{X_{12}}$, see
Divisors, Lemma \ref{divisors-lemma-strict-transform}.
Let $X_{12}'$ be the blowup of $X_{12}$ in
$p_1^{-1}\mathcal{I}_1 p_2^{-1}\mathcal{I}_2 \mathcal{O}_{X_{12}}$, see
Divisors, Lemma \ref{divisors-lemma-blowing-up-two-ideals}
for what this entails. We obtain in particular a commutative diagram
$$
\xymatrix{
X_{12}' \ar[d] \ar[r] & X_{12}^2 \ar[d] \\
X_{12}^1 \ar[r] & X_{12}
}
$$
where all the morphisms are $U$-admissible blowing ups.
Since $X_{12}^i \subset X_i^i$ is an open we may choose a $U$-admissible blowup
$X_i' \to X_i^i$ restricting to $X_{12}' \to X_{12}^i$, see
Divisors, Lemma \ref{divisors-lemma-extend-admissible-blowups}.
Then $X_{12}' \subset X_i'$ is an open subscheme and the diagram
$$
\xymatrix{
X_{12}' \ar[d] \ar[r] & X_i' \ar[d] \\
X_{12}^i \ar[r] & X_i^i
}
$$
is commutative with vertical arrows blowing ups and horizontal arrows
open immersions. Note that $X'_{12} \to X_1' \times_S X_2'$ is
an immersion and proper (use that $X'_{12} \to X_{12}$ is proper
and $X_{12} \to X_1 \times_S X_2$ is closed and $X_1' \times_S X_2' \to
X_1 \times_S X_2$ is separated and apply Morphisms, Lemma
\ref{morphisms-lemma-image-proper-scheme-closed}).
Thus $X'_{12} \to  X_1' \times_S X_2'$ is a closed immersion.
It follows that if we define $X$ by glueing $X_1'$ and $X_2'$
along the common open subscheme $X_{12}'$, then $X \to S$ is of finite type
and separated (Lemma \ref{lemma-check-separated}).
As compositions of $U$-admissible blowups are $U$-admissible blowups
(Divisors, Lemma \ref{divisors-lemma-composition-admissible-blowups})
the lemma is proved.
\end{proof}

\end{document}
